\documentclass[11pt]{article}
\newcommand{\LabNumber}{6}
\newcommand{\LabTitle}{Directives, Macros \& Checksum}
\newcommand{\LabTopic}{Assembler directives, AVR addressing modes, user-defined macros, and EEPROM data memory, consolidated into an Intel-HEX-style checksum integrity check -- the last assembly lab before the midterm}
\newcommand{\LabDuration}{3 hours (in-lab)}
% Shared preamble for Computer Programming lab sheets.
% Include with % Shared preamble for Computer Programming lab sheets.
% Include with % Shared preamble for Computer Programming lab sheets.
% Include with \input{common.tex} after \documentclass{article}.

\usepackage[a4paper,margin=2.2cm]{geometry}
\usepackage[T1]{fontenc}
\usepackage{lmodern}
\usepackage{microtype}
\usepackage{enumitem}
\usepackage{xcolor}
\usepackage{tcolorbox}
\tcbuselibrary{breakable,skins}
\usepackage{listings}
\usepackage{fancyhdr}
\usepackage{titlesec}
\usepackage{amsmath}
\usepackage{array}
\usepackage{booktabs}
\usepackage{hyperref}

\hypersetup{
  colorlinks=true,
  linkcolor=black,
  urlcolor=blue!60!black,
  pdfborder={0 0 0}
}

% ---------- Course identity (override per file if needed) ----------
\providecommand{\CourseName}{Microprocessors}
\providecommand{\CourseCode}{010153521}
\providecommand{\Semester}{Semester 1/2026}
\providecommand{\LabDuration}{3 hours (in-lab)}

% ---------- Colors ----------
\definecolor{accent}{HTML}{1F4E79}
\definecolor{accentlight}{HTML}{D9E2EC}
\definecolor{warn}{HTML}{B45309}
\definecolor{ok}{HTML}{166534}
\definecolor{codebg}{HTML}{F6F8FA}
\definecolor{codeframe}{HTML}{D0D7DE}
\definecolor{codekw}{HTML}{0550AE}
\definecolor{codestr}{HTML}{0A7B2C}
\definecolor{codecom}{HTML}{6E7781}

% ---------- Code / assembly listings ----------
\lstdefinestyle{c}{
  language=C,
  basicstyle=\ttfamily\small,
  keywordstyle=\color{codekw}\bfseries,
  stringstyle=\color{codestr},
  commentstyle=\color{codecom}\itshape,
  backgroundcolor=\color{codebg},
  frame=single,
  rulecolor=\color{codeframe},
  framerule=0.4pt,
  showstringspaces=false,
  breaklines=true,
  columns=fullflexible,
  keepspaces=true,
  tabsize=4,
  xleftmargin=8pt,
  xrightmargin=4pt,
  aboveskip=6pt,
  belowskip=6pt,
}
\lstdefinestyle{asm}{
  basicstyle=\ttfamily\small,
  commentstyle=\color{codecom}\itshape,
  backgroundcolor=\color{codebg},
  frame=single,
  rulecolor=\color{codeframe},
  framerule=0.4pt,
  showstringspaces=false,
  breaklines=true,
  columns=fullflexible,
  keepspaces=true,
  tabsize=4,
  xleftmargin=8pt,
  xrightmargin=4pt,
  aboveskip=6pt,
  belowskip=6pt,
}
\lstset{style=asm}

% ---------- Section styling ----------
\titleformat{\section}{\Large\bfseries\color{accent}}{\thesection.}{0.5em}{}
\titleformat{\subsection}{\large\bfseries\color{accent!85!black}}{\thesubsection}{0.5em}{}
\titlespacing*{\section}{0pt}{12pt}{6pt}
\titlespacing*{\subsection}{0pt}{8pt}{4pt}

% ---------- Page headers/footers ----------
\pagestyle{fancy}
\fancyhf{}
\fancyhead[L]{\small\CourseName\ifx\CourseCode\empty\else\ (\CourseCode)\fi}
\fancyhead[C]{\small\Semester}
\fancyhead[R]{\small\textbf{Lab \LabNumber:} \LabTitle}
\fancyfoot[C]{\small\thepage}
\renewcommand{\headrulewidth}{0.4pt}

% ---------- Title block ----------
\newcommand{\labheader}{%
  \begin{tcolorbox}[
    enhanced,
    colback=accentlight,
    colframe=accent,
    boxrule=0.5pt,
    arc=2pt,
    left=10pt, right=10pt, top=8pt, bottom=8pt
  ]
    {\Large\bfseries\color{accent} Lab \LabNumber: \LabTitle}\\[2pt]
    {\small \CourseName\ifx\CourseCode\empty\else\ (\CourseCode)\fi\quad\textbar\quad \Semester}\\[1pt]
    {\small \textbf{Duration:} \LabDuration\quad\textbar\quad \textbf{Topic:} \LabTopic}
  \end{tcolorbox}
  \vspace{2pt}
}

% ---------- Custom environments ----------
\newtcolorbox{summary}[1][Quick Reference]{
  enhanced,breakable,
  colback=accentlight!40,colframe=accent!70,
  title={\bfseries #1},
  fonttitle=\color{white},
  attach boxed title to top left={xshift=6pt,yshift=-3pt},
  boxed title style={colback=accent,boxrule=0pt,arc=2pt},
  arc=2pt,boxrule=0.4pt,
  top=12pt,
}

\newcounter{labex}[section]
\newtcolorbox{labexercise}[1]{
  enhanced,breakable,
  colback=white,colframe=accent,
  title={\bfseries In-Lab Exercise \refstepcounter{labex}\thelabex: #1},
  fonttitle=\color{white},
  attach boxed title to top left={xshift=6pt,yshift=-3pt},
  boxed title style={colback=accent,boxrule=0pt,arc=2pt},
  arc=2pt,boxrule=0.4pt,
  top=12pt,
}

\newcounter{traceex}
\newtcolorbox{trace}[1][]{
  enhanced,breakable,
  colback=white,colframe=warn,
  title={\bfseries Trace \& Predict \refstepcounter{traceex}\thetraceex\ifx\\#1\\\else: #1\fi},
  fonttitle=\color{white},
  attach boxed title to top left={xshift=6pt,yshift=-3pt},
  boxed title style={colback=warn,boxrule=0pt,arc=2pt},
  arc=2pt,boxrule=0.4pt,
  top=12pt,
}

\newcounter{bugex}
\newtcolorbox{bug}[1][]{
  enhanced,breakable,
  colback=white,colframe=warn,
  title={\bfseries Find the Bug \refstepcounter{bugex}\thebugex\ifx\\#1\\\else: #1\fi},
  fonttitle=\color{white},
  attach boxed title to top left={xshift=6pt,yshift=-3pt},
  boxed title style={colback=warn,boxrule=0pt,arc=2pt},
  arc=2pt,boxrule=0.4pt,
  top=12pt,
}

\newcounter{writeex}
\newtcolorbox{writecode}[1][]{
  enhanced,breakable,
  colback=white,colframe=warn,
  title={\bfseries Write Code on Paper \refstepcounter{writeex}\thewriteex\ifx\\#1\\\else: #1\fi},
  fonttitle=\color{white},
  attach boxed title to top left={xshift=6pt,yshift=-3pt},
  boxed title style={colback=warn,boxrule=0pt,arc=2pt},
  arc=2pt,boxrule=0.4pt,
  top=12pt,
}

% ---------- AI Collaboration box ----------
\definecolor{aipurple}{HTML}{4C1D95}
\definecolor{aipurplelight}{HTML}{EDE9FE}

\newtcolorbox{aicollab}[1][AI Collaboration]{
  enhanced,breakable,
  colback=aipurplelight,colframe=aipurple!70,
  title={\bfseries #1},
  fonttitle=\color{white},
  attach boxed title to top left={xshift=6pt,yshift=-3pt},
  boxed title style={colback=aipurple,boxrule=0pt,arc=2pt},
  arc=2pt,boxrule=0.4pt,
  top=12pt,
}

\newcounter{tryex}
\newtcolorbox{tryit}[1]{
  enhanced,breakable,
  colback=ok!6,colframe=ok!70!black,
  title={\bfseries Try It \refstepcounter{tryex}\thetryex: #1},
  fonttitle=\color{white},
  attach boxed title to top left={xshift=6pt,yshift=-3pt},
  boxed title style={colback=ok!70!black,boxrule=0pt,arc=2pt},
  arc=2pt,boxrule=0.4pt,
  top=12pt,
}

% ---------- Drawing exercise ----------
\definecolor{drawteal}{HTML}{0D7377}
\definecolor{drawteallight}{HTML}{D4F1F4}

\newcounter{drawex}
\newtcolorbox{drawex}[1]{
  enhanced,breakable,
  colback=drawteallight,colframe=drawteal!70,
  title={\bfseries Draw on Paper \refstepcounter{drawex}\thedrawex: #1},
  fonttitle=\color{white},
  attach boxed title to top left={xshift=6pt,yshift=-3pt},
  boxed title style={colback=drawteal,boxrule=0pt,arc=2pt},
  arc=2pt,boxrule=0.4pt,
  top=12pt,
}


\newcommand{\takehomebanner}{%
  \vspace{6pt}
  \begin{tcolorbox}[
    colback=warn!10,colframe=warn,boxrule=0.5pt,arc=2pt,
    left=10pt,right=10pt,top=6pt,bottom=6pt
  ]
    \textbf{Take-Home (Paper Work).}\;
    Solve the following on paper without using a computer or compiler.
    Hand-write your answers; submit at the start of the next lab.
    Goal: prepare you to dry-code during the written exam.
  \end{tcolorbox}
  \vspace{4pt}
}
 after \documentclass{article}.

\usepackage[a4paper,margin=2.2cm]{geometry}
\usepackage[T1]{fontenc}
\usepackage{lmodern}
\usepackage{microtype}
\usepackage{enumitem}
\usepackage{xcolor}
\usepackage{tcolorbox}
\tcbuselibrary{breakable,skins}
\usepackage{listings}
\usepackage{fancyhdr}
\usepackage{titlesec}
\usepackage{amsmath}
\usepackage{array}
\usepackage{booktabs}
\usepackage{hyperref}

\hypersetup{
  colorlinks=true,
  linkcolor=black,
  urlcolor=blue!60!black,
  pdfborder={0 0 0}
}

% ---------- Course identity (override per file if needed) ----------
\providecommand{\CourseName}{Microprocessors}
\providecommand{\CourseCode}{010153521}
\providecommand{\Semester}{Semester 1/2026}
\providecommand{\LabDuration}{3 hours (in-lab)}

% ---------- Colors ----------
\definecolor{accent}{HTML}{1F4E79}
\definecolor{accentlight}{HTML}{D9E2EC}
\definecolor{warn}{HTML}{B45309}
\definecolor{ok}{HTML}{166534}
\definecolor{codebg}{HTML}{F6F8FA}
\definecolor{codeframe}{HTML}{D0D7DE}
\definecolor{codekw}{HTML}{0550AE}
\definecolor{codestr}{HTML}{0A7B2C}
\definecolor{codecom}{HTML}{6E7781}

% ---------- Code / assembly listings ----------
\lstdefinestyle{c}{
  language=C,
  basicstyle=\ttfamily\small,
  keywordstyle=\color{codekw}\bfseries,
  stringstyle=\color{codestr},
  commentstyle=\color{codecom}\itshape,
  backgroundcolor=\color{codebg},
  frame=single,
  rulecolor=\color{codeframe},
  framerule=0.4pt,
  showstringspaces=false,
  breaklines=true,
  columns=fullflexible,
  keepspaces=true,
  tabsize=4,
  xleftmargin=8pt,
  xrightmargin=4pt,
  aboveskip=6pt,
  belowskip=6pt,
}
\lstdefinestyle{asm}{
  basicstyle=\ttfamily\small,
  commentstyle=\color{codecom}\itshape,
  backgroundcolor=\color{codebg},
  frame=single,
  rulecolor=\color{codeframe},
  framerule=0.4pt,
  showstringspaces=false,
  breaklines=true,
  columns=fullflexible,
  keepspaces=true,
  tabsize=4,
  xleftmargin=8pt,
  xrightmargin=4pt,
  aboveskip=6pt,
  belowskip=6pt,
}
\lstset{style=asm}

% ---------- Section styling ----------
\titleformat{\section}{\Large\bfseries\color{accent}}{\thesection.}{0.5em}{}
\titleformat{\subsection}{\large\bfseries\color{accent!85!black}}{\thesubsection}{0.5em}{}
\titlespacing*{\section}{0pt}{12pt}{6pt}
\titlespacing*{\subsection}{0pt}{8pt}{4pt}

% ---------- Page headers/footers ----------
\pagestyle{fancy}
\fancyhf{}
\fancyhead[L]{\small\CourseName\ifx\CourseCode\empty\else\ (\CourseCode)\fi}
\fancyhead[C]{\small\Semester}
\fancyhead[R]{\small\textbf{Lab \LabNumber:} \LabTitle}
\fancyfoot[C]{\small\thepage}
\renewcommand{\headrulewidth}{0.4pt}

% ---------- Title block ----------
\newcommand{\labheader}{%
  \begin{tcolorbox}[
    enhanced,
    colback=accentlight,
    colframe=accent,
    boxrule=0.5pt,
    arc=2pt,
    left=10pt, right=10pt, top=8pt, bottom=8pt
  ]
    {\Large\bfseries\color{accent} Lab \LabNumber: \LabTitle}\\[2pt]
    {\small \CourseName\ifx\CourseCode\empty\else\ (\CourseCode)\fi\quad\textbar\quad \Semester}\\[1pt]
    {\small \textbf{Duration:} \LabDuration\quad\textbar\quad \textbf{Topic:} \LabTopic}
  \end{tcolorbox}
  \vspace{2pt}
}

% ---------- Custom environments ----------
\newtcolorbox{summary}[1][Quick Reference]{
  enhanced,breakable,
  colback=accentlight!40,colframe=accent!70,
  title={\bfseries #1},
  fonttitle=\color{white},
  attach boxed title to top left={xshift=6pt,yshift=-3pt},
  boxed title style={colback=accent,boxrule=0pt,arc=2pt},
  arc=2pt,boxrule=0.4pt,
  top=12pt,
}

\newcounter{labex}[section]
\newtcolorbox{labexercise}[1]{
  enhanced,breakable,
  colback=white,colframe=accent,
  title={\bfseries In-Lab Exercise \refstepcounter{labex}\thelabex: #1},
  fonttitle=\color{white},
  attach boxed title to top left={xshift=6pt,yshift=-3pt},
  boxed title style={colback=accent,boxrule=0pt,arc=2pt},
  arc=2pt,boxrule=0.4pt,
  top=12pt,
}

\newcounter{traceex}
\newtcolorbox{trace}[1][]{
  enhanced,breakable,
  colback=white,colframe=warn,
  title={\bfseries Trace \& Predict \refstepcounter{traceex}\thetraceex\ifx\\#1\\\else: #1\fi},
  fonttitle=\color{white},
  attach boxed title to top left={xshift=6pt,yshift=-3pt},
  boxed title style={colback=warn,boxrule=0pt,arc=2pt},
  arc=2pt,boxrule=0.4pt,
  top=12pt,
}

\newcounter{bugex}
\newtcolorbox{bug}[1][]{
  enhanced,breakable,
  colback=white,colframe=warn,
  title={\bfseries Find the Bug \refstepcounter{bugex}\thebugex\ifx\\#1\\\else: #1\fi},
  fonttitle=\color{white},
  attach boxed title to top left={xshift=6pt,yshift=-3pt},
  boxed title style={colback=warn,boxrule=0pt,arc=2pt},
  arc=2pt,boxrule=0.4pt,
  top=12pt,
}

\newcounter{writeex}
\newtcolorbox{writecode}[1][]{
  enhanced,breakable,
  colback=white,colframe=warn,
  title={\bfseries Write Code on Paper \refstepcounter{writeex}\thewriteex\ifx\\#1\\\else: #1\fi},
  fonttitle=\color{white},
  attach boxed title to top left={xshift=6pt,yshift=-3pt},
  boxed title style={colback=warn,boxrule=0pt,arc=2pt},
  arc=2pt,boxrule=0.4pt,
  top=12pt,
}

% ---------- AI Collaboration box ----------
\definecolor{aipurple}{HTML}{4C1D95}
\definecolor{aipurplelight}{HTML}{EDE9FE}

\newtcolorbox{aicollab}[1][AI Collaboration]{
  enhanced,breakable,
  colback=aipurplelight,colframe=aipurple!70,
  title={\bfseries #1},
  fonttitle=\color{white},
  attach boxed title to top left={xshift=6pt,yshift=-3pt},
  boxed title style={colback=aipurple,boxrule=0pt,arc=2pt},
  arc=2pt,boxrule=0.4pt,
  top=12pt,
}

\newcounter{tryex}
\newtcolorbox{tryit}[1]{
  enhanced,breakable,
  colback=ok!6,colframe=ok!70!black,
  title={\bfseries Try It \refstepcounter{tryex}\thetryex: #1},
  fonttitle=\color{white},
  attach boxed title to top left={xshift=6pt,yshift=-3pt},
  boxed title style={colback=ok!70!black,boxrule=0pt,arc=2pt},
  arc=2pt,boxrule=0.4pt,
  top=12pt,
}

% ---------- Drawing exercise ----------
\definecolor{drawteal}{HTML}{0D7377}
\definecolor{drawteallight}{HTML}{D4F1F4}

\newcounter{drawex}
\newtcolorbox{drawex}[1]{
  enhanced,breakable,
  colback=drawteallight,colframe=drawteal!70,
  title={\bfseries Draw on Paper \refstepcounter{drawex}\thedrawex: #1},
  fonttitle=\color{white},
  attach boxed title to top left={xshift=6pt,yshift=-3pt},
  boxed title style={colback=drawteal,boxrule=0pt,arc=2pt},
  arc=2pt,boxrule=0.4pt,
  top=12pt,
}


\newcommand{\takehomebanner}{%
  \vspace{6pt}
  \begin{tcolorbox}[
    colback=warn!10,colframe=warn,boxrule=0.5pt,arc=2pt,
    left=10pt,right=10pt,top=6pt,bottom=6pt
  ]
    \textbf{Take-Home (Paper Work).}\;
    Solve the following on paper without using a computer or compiler.
    Hand-write your answers; submit at the start of the next lab.
    Goal: prepare you to dry-code during the written exam.
  \end{tcolorbox}
  \vspace{4pt}
}
 after \documentclass{article}.

\usepackage[a4paper,margin=2.2cm]{geometry}
\usepackage[T1]{fontenc}
\usepackage{lmodern}
\usepackage{microtype}
\usepackage{enumitem}
\usepackage{xcolor}
\usepackage{tcolorbox}
\tcbuselibrary{breakable,skins}
\usepackage{listings}
\usepackage{fancyhdr}
\usepackage{titlesec}
\usepackage{amsmath}
\usepackage{array}
\usepackage{booktabs}
\usepackage{hyperref}

\hypersetup{
  colorlinks=true,
  linkcolor=black,
  urlcolor=blue!60!black,
  pdfborder={0 0 0}
}

% ---------- Course identity (override per file if needed) ----------
\providecommand{\CourseName}{Microprocessors}
\providecommand{\CourseCode}{010153521}
\providecommand{\Semester}{Semester 1/2026}
\providecommand{\LabDuration}{3 hours (in-lab)}

% ---------- Colors ----------
\definecolor{accent}{HTML}{1F4E79}
\definecolor{accentlight}{HTML}{D9E2EC}
\definecolor{warn}{HTML}{B45309}
\definecolor{ok}{HTML}{166534}
\definecolor{codebg}{HTML}{F6F8FA}
\definecolor{codeframe}{HTML}{D0D7DE}
\definecolor{codekw}{HTML}{0550AE}
\definecolor{codestr}{HTML}{0A7B2C}
\definecolor{codecom}{HTML}{6E7781}

% ---------- Code / assembly listings ----------
\lstdefinestyle{c}{
  language=C,
  basicstyle=\ttfamily\small,
  keywordstyle=\color{codekw}\bfseries,
  stringstyle=\color{codestr},
  commentstyle=\color{codecom}\itshape,
  backgroundcolor=\color{codebg},
  frame=single,
  rulecolor=\color{codeframe},
  framerule=0.4pt,
  showstringspaces=false,
  breaklines=true,
  columns=fullflexible,
  keepspaces=true,
  tabsize=4,
  xleftmargin=8pt,
  xrightmargin=4pt,
  aboveskip=6pt,
  belowskip=6pt,
}
\lstdefinestyle{asm}{
  basicstyle=\ttfamily\small,
  commentstyle=\color{codecom}\itshape,
  backgroundcolor=\color{codebg},
  frame=single,
  rulecolor=\color{codeframe},
  framerule=0.4pt,
  showstringspaces=false,
  breaklines=true,
  columns=fullflexible,
  keepspaces=true,
  tabsize=4,
  xleftmargin=8pt,
  xrightmargin=4pt,
  aboveskip=6pt,
  belowskip=6pt,
}
\lstset{style=asm}

% ---------- Section styling ----------
\titleformat{\section}{\Large\bfseries\color{accent}}{\thesection.}{0.5em}{}
\titleformat{\subsection}{\large\bfseries\color{accent!85!black}}{\thesubsection}{0.5em}{}
\titlespacing*{\section}{0pt}{12pt}{6pt}
\titlespacing*{\subsection}{0pt}{8pt}{4pt}

% ---------- Page headers/footers ----------
\pagestyle{fancy}
\fancyhf{}
\fancyhead[L]{\small\CourseName\ifx\CourseCode\empty\else\ (\CourseCode)\fi}
\fancyhead[C]{\small\Semester}
\fancyhead[R]{\small\textbf{Lab \LabNumber:} \LabTitle}
\fancyfoot[C]{\small\thepage}
\renewcommand{\headrulewidth}{0.4pt}

% ---------- Title block ----------
\newcommand{\labheader}{%
  \begin{tcolorbox}[
    enhanced,
    colback=accentlight,
    colframe=accent,
    boxrule=0.5pt,
    arc=2pt,
    left=10pt, right=10pt, top=8pt, bottom=8pt
  ]
    {\Large\bfseries\color{accent} Lab \LabNumber: \LabTitle}\\[2pt]
    {\small \CourseName\ifx\CourseCode\empty\else\ (\CourseCode)\fi\quad\textbar\quad \Semester}\\[1pt]
    {\small \textbf{Duration:} \LabDuration\quad\textbar\quad \textbf{Topic:} \LabTopic}
  \end{tcolorbox}
  \vspace{2pt}
}

% ---------- Custom environments ----------
\newtcolorbox{summary}[1][Quick Reference]{
  enhanced,breakable,
  colback=accentlight!40,colframe=accent!70,
  title={\bfseries #1},
  fonttitle=\color{white},
  attach boxed title to top left={xshift=6pt,yshift=-3pt},
  boxed title style={colback=accent,boxrule=0pt,arc=2pt},
  arc=2pt,boxrule=0.4pt,
  top=12pt,
}

\newcounter{labex}[section]
\newtcolorbox{labexercise}[1]{
  enhanced,breakable,
  colback=white,colframe=accent,
  title={\bfseries In-Lab Exercise \refstepcounter{labex}\thelabex: #1},
  fonttitle=\color{white},
  attach boxed title to top left={xshift=6pt,yshift=-3pt},
  boxed title style={colback=accent,boxrule=0pt,arc=2pt},
  arc=2pt,boxrule=0.4pt,
  top=12pt,
}

\newcounter{traceex}
\newtcolorbox{trace}[1][]{
  enhanced,breakable,
  colback=white,colframe=warn,
  title={\bfseries Trace \& Predict \refstepcounter{traceex}\thetraceex\ifx\\#1\\\else: #1\fi},
  fonttitle=\color{white},
  attach boxed title to top left={xshift=6pt,yshift=-3pt},
  boxed title style={colback=warn,boxrule=0pt,arc=2pt},
  arc=2pt,boxrule=0.4pt,
  top=12pt,
}

\newcounter{bugex}
\newtcolorbox{bug}[1][]{
  enhanced,breakable,
  colback=white,colframe=warn,
  title={\bfseries Find the Bug \refstepcounter{bugex}\thebugex\ifx\\#1\\\else: #1\fi},
  fonttitle=\color{white},
  attach boxed title to top left={xshift=6pt,yshift=-3pt},
  boxed title style={colback=warn,boxrule=0pt,arc=2pt},
  arc=2pt,boxrule=0.4pt,
  top=12pt,
}

\newcounter{writeex}
\newtcolorbox{writecode}[1][]{
  enhanced,breakable,
  colback=white,colframe=warn,
  title={\bfseries Write Code on Paper \refstepcounter{writeex}\thewriteex\ifx\\#1\\\else: #1\fi},
  fonttitle=\color{white},
  attach boxed title to top left={xshift=6pt,yshift=-3pt},
  boxed title style={colback=warn,boxrule=0pt,arc=2pt},
  arc=2pt,boxrule=0.4pt,
  top=12pt,
}

% ---------- AI Collaboration box ----------
\definecolor{aipurple}{HTML}{4C1D95}
\definecolor{aipurplelight}{HTML}{EDE9FE}

\newtcolorbox{aicollab}[1][AI Collaboration]{
  enhanced,breakable,
  colback=aipurplelight,colframe=aipurple!70,
  title={\bfseries #1},
  fonttitle=\color{white},
  attach boxed title to top left={xshift=6pt,yshift=-3pt},
  boxed title style={colback=aipurple,boxrule=0pt,arc=2pt},
  arc=2pt,boxrule=0.4pt,
  top=12pt,
}

\newcounter{tryex}
\newtcolorbox{tryit}[1]{
  enhanced,breakable,
  colback=ok!6,colframe=ok!70!black,
  title={\bfseries Try It \refstepcounter{tryex}\thetryex: #1},
  fonttitle=\color{white},
  attach boxed title to top left={xshift=6pt,yshift=-3pt},
  boxed title style={colback=ok!70!black,boxrule=0pt,arc=2pt},
  arc=2pt,boxrule=0.4pt,
  top=12pt,
}

% ---------- Drawing exercise ----------
\definecolor{drawteal}{HTML}{0D7377}
\definecolor{drawteallight}{HTML}{D4F1F4}

\newcounter{drawex}
\newtcolorbox{drawex}[1]{
  enhanced,breakable,
  colback=drawteallight,colframe=drawteal!70,
  title={\bfseries Draw on Paper \refstepcounter{drawex}\thedrawex: #1},
  fonttitle=\color{white},
  attach boxed title to top left={xshift=6pt,yshift=-3pt},
  boxed title style={colback=drawteal,boxrule=0pt,arc=2pt},
  arc=2pt,boxrule=0.4pt,
  top=12pt,
}


\newcommand{\takehomebanner}{%
  \vspace{6pt}
  \begin{tcolorbox}[
    colback=warn!10,colframe=warn,boxrule=0.5pt,arc=2pt,
    left=10pt,right=10pt,top=6pt,bottom=6pt
  ]
    \textbf{Take-Home (Paper Work).}\;
    Solve the following on paper without using a computer or compiler.
    Hand-write your answers; submit at the start of the next lab.
    Goal: prepare you to dry-code during the written exam.
  \end{tcolorbox}
  \vspace{4pt}
}

\usepackage{amssymb}

\begin{document}
\labheader

\section*{Objectives}
\begin{itemize}[leftmargin=*,nosep]
  \item Use assembler \textbf{directives} to organize a program across Flash code
        (\texttt{.CSEG}), a Flash data table (\texttt{.DB}), and EEPROM data
        (\texttt{.ESEG}) -- not just \texttt{.INCLUDE} and \texttt{.ORG}, which
        you've already used since Lab 1.
  \item Contrast AVR's \textbf{addressing modes} side by side: immediate,
        register-direct, and indirect-with-post-increment (the \texttt{Z} register
        with \texttt{LPM}) to walk a table in Flash.
  \item Write and use your own \textbf{macro} (\texttt{.MACRO}/\texttt{.ENDM}) to
        eliminate repeated code -- and learn the one gotcha every AVR macro writer
        hits eventually.
  \item Read persistent \textbf{EEPROM} memory at runtime and compare it against a
        freshly computed \textbf{checksum} -- the same two's-complement scheme
        every Intel HEX file you've uploaded since Lab 3 already relies on.
\end{itemize}

\begin{summary}[Quick Reference: Assembler Directives You'll Use]
\begin{center}
\begin{tabular}{@{}p{2.3cm}p{6.0cm}p{5.6cm}@{}}
\toprule
\textbf{Directive} & \textbf{Purpose} & \textbf{Example} \\
\midrule
\texttt{.EQU} & Name a constant -- uses no memory, just text substitution at
assemble time & \texttt{.EQU TABLE\_LEN = 8} \\
\texttt{.CSEG} & Begin/continue the Flash \emph{code} segment (also holds
\texttt{.DB} tables read via \texttt{LPM}) & \texttt{.CSEG} \\
\texttt{.DB} & Define byte(s) at the current address in the current segment &
\texttt{TABLE: .DB 8,4,7,5,6,1,2,3} \\
\texttt{.ESEG} & Begin/continue the \textbf{EEPROM} segment -- a separate address
space from Flash and SRAM & \texttt{.ESEG} \\
\texttt{.ORG} & Set the address of whatever follows, within the current segment &
\texttt{.ORG 0x0000} \\
\texttt{.MACRO} / \texttt{.ENDM} & Define a reusable named block of instructions,
expanded inline at every invocation & see the macro box below \\
\bottomrule
\end{tabular}
\end{center}
Every directive here is scoped to \emph{assembly time} -- none of them exist once
your program is running. The three segments (\texttt{.CSEG}/Flash,
\texttt{.ESEG}/EEPROM, and SRAM's \texttt{.DSEG}, not used in this lab) are
physically separate memories with separate address spaces: address \texttt{0x0000}
in Flash and address \texttt{0x0000} in EEPROM are two completely different
storage cells.
\end{summary}

\begin{summary}[Quick Reference: Addressing Modes Side by Side]
\begin{center}
\begin{tabular}{@{}p{3.6cm}p{3.6cm}p{6.7cm}@{}}
\toprule
\textbf{Mode} & \textbf{Example} & \textbf{What it means} \\
\midrule
Immediate & \texttt{LDI R16,0x24} & The operand \emph{is} the value \texttt{0x24},
encoded directly in the instruction word \\
Register (direct) & \texttt{ADD R17,R16} & The operand is whatever a register
currently holds \\
Direct I/O & \texttt{OUT PORTB,R16} & The operand is a fixed I/O register address \\
Indirect, post-increment & \texttt{LPM R16,Z+} & The address to read is held
\emph{in} the \texttt{Z} register (a pointer, not a value); read Flash at that
address into \texttt{R16}, then increment \texttt{Z} \\
\bottomrule
\end{tabular}
\end{center}
Every AVR instruction you've written through Lab 5 used one of the first three
modes. Indirect addressing is the new one here: instead of the instruction naming
a fixed address, a register \emph{holds} the address -- which is exactly what lets
a fixed-size loop walk a table of any length without one hardcoded instruction per
byte.
\end{summary}

\begin{summary}[Quick Reference: Defining and Using a Macro]
\begin{lstlisting}
.MACRO CHECKSUM_STEP
    LPM   R16, Z+      ;read next table byte, advance Z
    ADD   R17, R16     ;add it into the running sum
.ENDM

;invoke it just like any instruction, as many times as you like:
CHECKSUM_STEP
CHECKSUM_STEP
\end{lstlisting}
A macro's body is \emph{inserted fresh} into your code at every invocation --
text substitution at assemble time, not a runtime call like \texttt{CALL}/\texttt{RET}
(no stack use, no return address, but also no code-size savings: 10 invocations
means 10 copies of the body in Flash). \texttt{CHECKSUM\_STEP} above is safe to
invoke any number of times because it contains no label of its own. A macro that
\emph{does} contain an internal label will fail to assemble the second time it's
invoked in the same file -- that label would then be defined twice (see Find the
Bug~1).
\end{summary}

\begin{summary}[Quick Reference: EEPROM Memory]
EEPROM is a third, separate memory: unlike SRAM it \textbf{survives power loss},
but writes are far slower ($\approx$3.3~ms/byte) and rated for only about
\textbf{100,000 write cycles per byte} on the ATmega328P -- reserve it for values
that must persist, not for anything rewritten every loop pass. Reserve space with
\texttt{.ESEG} exactly like a Flash table, but reading it back at runtime needs the
EEPROM control registers, not \texttt{LPM}:
\begin{lstlisting}
.MACRO EEPROM_READ           ;always reads EEPROM address 0x0000
WAIT_EE:
    SBIC  EECR, EEPE         ;wait for any earlier write to finish
    RJMP  WAIT_EE
    LDI   R16, 0
    OUT   EEARH, R16
    OUT   EEARL, R16         ;address = 0x0000
    SBI   EECR, EERE         ;trigger the read
    IN    @0, EEDR           ;@0 = the macro's parameter (destination register)
.ENDM
\end{lstlisting}
(\texttt{EEPE}/\texttt{EERE} are \texttt{m328Pdef.inc}'s names for what older
textbooks and chips call \texttt{EEWE}/\texttt{EERE} -- same register bits, newer
names, exactly the kind of naming difference Lab 3 already warned you about.)
Studio assembles a separate \texttt{.eep} file alongside your \texttt{.hex};
upload it with \texttt{avrdude}'s EEPROM memory type -- the same \texttt{-U}
syntax as Lab 3's flash write, just a different memory name:
\begin{lstlisting}[language=bash]
avrdude -c arduino -p atmega328p -P <port> -b 115200 -U eeprom:w:yourfile.eep:i
\end{lstlisting}
\end{summary}

\begin{summary}[Quick Reference: The Two's-Complement Checksum]
\[
\text{sum} = (\text{byte}_1 + \text{byte}_2 + \cdots + \text{byte}_N) \bmod 256,
\qquad
\text{checksum} = (256 - \text{sum}) \bmod 256
\]
\texttt{checksum} is simply the two's complement of \texttt{sum} (one
\texttt{NEG} instruction). This definition guarantees
$\text{sum} + \text{checksum} \equiv 0 \pmod{256}$ always -- so \emph{verifying}
is just ``add every byte, including the checksum byte itself; the total's low
byte must come out exactly \texttt{0}.'' This is the \emph{exact} scheme Intel HEX
records use: \texttt{avrdude} has been checking a checksum exactly like this on
every record of every \texttt{.hex} upload you've done since Lab 3.
\end{summary}

\begin{summary}[Quick Reference: Your Individual Data Table]
Take the \textbf{last 8 digits} of your student ID as 8 individual byte values
(\texttt{0-9} each) -- this is your \texttt{TABLE}.
\textbf{Worked example} (a made-up ID, do not reuse): ID \texttt{1029384756123}
$\to$ last 8 digits \texttt{84756123} $\to$ table bytes
\texttt{8,4,7,5,6,1,2,3}.
\[
\text{sum} = 8{+}4{+}7{+}5{+}6{+}1{+}2{+}3 = 36\ (\texttt{0x24}), \qquad
\text{checksum} = 256 - 36 = 220\ (\texttt{0xDC})
\]
Check: $36 + 220 = 256$, low byte $= 0$. \checkmark\ Compute \emph{your own} table
and checksum now, by hand, in Exercise 1 -- before touching the keyboard.
\end{summary}

\section*{Bill of Materials}
No new parts this lab -- just the \textbf{Arduino Uno R3} from previous labs. This
lab reuses the Uno's \textbf{on-board LED} (\texttt{PB5}/D13) for pass/fail
indication; no breadboard circuit is required.

\section*{Quick Experiments (Try It)}

\begin{tryit}{Your first macro}
Define \texttt{INIT\_STACK} as a macro (the four-line stack-pointer
initialization you've hand-typed at the top of every program since Lab 3):
\begin{lstlisting}
.MACRO INIT_STACK
    LDI   R16, HIGH(RAMEND)
    OUT   SPH, R16
    LDI   R16, LOW(RAMEND)
    OUT   SPL, R16
.ENDM
\end{lstlisting}
Invoke it once at the top of \texttt{MAIN}, build, and open the \texttt{.lst}
file -- confirm the single line \texttt{INIT\_STACK} expanded into the same four
instructions you'd have typed by hand.
\end{tryit}

\begin{tryit}{One indirect read, by itself}
Load \texttt{Z} with your \texttt{TABLE}'s address (\texttt{LDI ZL,LOW(TABLE*2)}
/ \texttt{LDI ZH,HIGH(TABLE*2)} -- Flash byte addresses are double the word
address the label gives you, hence the \texttt{*2}), execute a single
\texttt{LPM R16,Z+} by itself, and confirm \texttt{R16} holds your table's
\emph{first} byte (write it to \texttt{PORTB} and check in the simulator, or on a
logic probe) before writing any loop.
\end{tryit}

\section*{In-Lab Exercises (target: 3 hours)}

\begin{labexercise}{Derive your table and checksum by hand \hfill {\small \itshape (25 min)}}
Using the ``Your Individual Data Table'' box above, write out your 8 table bytes
and compute \texttt{sum} and \texttt{checksum} by hand in your notebook, showing
every step. Double-check by adding \texttt{sum + checksum} yourself and confirming
the low byte is exactly \texttt{0} -- this is the same value your program will
compute in Exercise 3, so an arithmetic slip here will cost you later.
\end{labexercise}

\begin{labexercise}{Program structure with directives \hfill {\small \itshape (35 min)}}
Set up your Microchip Studio project with:
\begin{itemize}[nosep,leftmargin=*]
  \item A \texttt{.CSEG} table (\texttt{TABLE: .DB} your 8 bytes from Exercise 1)
        and \texttt{.EQU TABLE\_LEN = 8}.
  \item An \texttt{.ESEG} reservation (\texttt{REF\_CHECKSUM: .DB} your
        hand-computed checksum from Exercise 1).
  \item Your \texttt{INIT\_STACK} macro from Try It 1, invoked once in
        \texttt{MAIN}.
\end{itemize}
Build the project -- Studio produces \emph{both} a \texttt{.hex} and a
\texttt{.eep} file. Upload \textbf{both} with \texttt{avrdude} (flash exactly as
in Lab 3, EEPROM using the Quick Reference box's \texttt{eeprom} memory type). Do
\emph{not} write any checksum-verification logic yet -- this exercise is only
about getting the segments, directives, and both uploads right.
\end{labexercise}

\begin{labexercise}{Runtime checksum and verification \hfill {\small \itshape (85 min)}}
Write your own \texttt{CHECKSUM\_STEP} macro (following the Quick Reference
pattern) and use it inside a loop of \texttt{TABLE\_LEN} iterations to compute the
actual checksum over your Flash \texttt{TABLE} at runtime, using \texttt{Z} +
\texttt{LPM} indirect addressing. Use the \emph{given} \texttt{EEPROM\_READ}
macro (Quick Reference box above) to read \texttt{REF\_CHECKSUM} back from
EEPROM. Compare the two values (\texttt{CP}); if they match, light the on-board
LED (\texttt{PB5}) steady; if they don't, leave it off.
\begin{itemize}[nosep,leftmargin=*]
  \item \textbf{Trace/predict first, then verify:} before running it, predict
        what your program's internally computed \texttt{sum + checksum} should
        equal, and what the LED should show. Then build, upload, and confirm.
\end{itemize}
\end{labexercise}

\begin{labexercise}{Break it on purpose \hfill {\small \itshape (30 min)}}
Deliberately change \emph{one} byte in your Flash \texttt{TABLE}, re-assemble,
and re-upload \textbf{flash only} (leave the EEPROM reference untouched). Confirm
the LED now correctly reports failure. Change the byte back, re-upload, and
confirm the LED reports success again.
\begin{itemize}[nosep,leftmargin=*]
  \item In your notebook, explain in one or two sentences why this exact
        mechanism -- a stored reference checksum compared against a freshly
        computed one -- is how real bootloaders and firmware updates detect a
        corrupted transfer.
\end{itemize}
\end{labexercise}

\begin{aicollab}
\textbf{Try these prompts with an AI tool:}
\begin{itemize}[leftmargin=*,nosep]
  \item ``What's the actual difference between an AVR macro
        (\texttt{.MACRO}/\texttt{.ENDM}) and a subroutine
        (\texttt{CALL}/\texttt{RET}) -- when would I choose one over the other?''
  \item ``Why does a fixed internal label inside an AVR macro cause a `duplicate
        label' error the second time the macro is invoked, and how do real AVR
        programs work around it?''
  \item ``What is EEPROM's typical write endurance on an ATmega328P, and why does
        that affect what you should and shouldn't store there?''
\end{itemize}

\medskip\textbf{Critical check -- verify, don't just accept:}
If your Exercise 3 LED shows failure on the first try, verify your hand-computed
checksum (Exercise 1) against the runtime value byte by byte before assuming your
code is wrong -- a surprising fraction of first failures turn out to be an
arithmetic slip in the \emph{hand} calculation, not a bug in the program.

\medskip\textbf{Know the limit:}
AI tools frequently use ``checksum'' and ``CRC'' interchangeably -- they are
\emph{not} the same algorithm and do not catch the same errors. A simple
sum-based checksum like this lab's misses several common multi-bit corruption
patterns that a CRC would catch; don't let an AI's phrasing imply your checksum
has CRC-level protection when asking it for help.
\end{aicollab}

\takehomebanner

\begin{trace}[Compute a checksum by hand]
For the 6-byte table \texttt{0x1A, 0x03, 0xFF, 0x40, 0x02, 0x1C}, compute
\texttt{sum} and the two's-complement \texttt{checksum} by hand, showing your
work. Verify by confirming \texttt{sum + checksum} has a low byte of \texttt{0}.
\end{trace}

\begin{bug}[A macro invoked twice]
A student wrote the \texttt{EEPROM\_READ} macro exactly as given in the Quick
Reference box (with its internal \texttt{WAIT\_EE} label) and invoked it
\emph{twice} in the same program -- once to read a reference checksum, and once
to read a second stored calibration byte at a different address. Predict the
exact category of error the assembler reports, and explain -- in terms of what a
macro invocation actually does to your source code -- why it happens on the
\emph{second} invocation specifically, not the first.
\end{bug}

\begin{writecode}[A parameterized macro]
On paper, without a computer, simulator, or AI tool, write a macro
\texttt{TOGGLE\_LED} that takes one parameter (a bit number, \texttt{0-7}) and
toggles that bit of \texttt{PORTB} -- \emph{without} first reading the pin to
figure out which state it's currently in. (Hint: re-read the ATmega328P port
section on what writing a \texttt{1} to a \texttt{PINx} bit does.)
\end{writecode}

\begin{drawex}{Draw your program's memory map}
Draw a memory map of your assembled Exercise 3 program: separate boxes for Flash
(showing \texttt{MAIN}'s code and \texttt{TABLE}, with its address) and EEPROM
(showing \texttt{REF\_CHECKSUM}, with its address) as two \emph{physically
separate} address spaces. For each box, label which instruction and which
addressing mode your program uses to reach it (e.g.\ ``\texttt{LPM Z+}, indirect,
reads \texttt{TABLE}'').
\end{drawex}

\section*{Reflection}

This is the last assembly lab before the midterm -- after this, the course moves
to C. In one short paragraph, look ahead: which of today's five concepts
(directives, addressing modes, macros, EEPROM, checksums) do you expect to become
\emph{invisible} in C, handled automatically by the compiler, and which will you
still control directly yourself, just through different syntax? Consider
specifically: assembler segments (\texttt{.CSEG}/\texttt{.ESEG}) versus a
compiler's memory sections, indirect addressing versus C pointers, \texttt{.MACRO}
versus \texttt{\#define}, and EEPROM access via raw registers versus a library
call. Is there anything on today's list that C \emph{doesn't} make easier?

\end{document}
